% Общий преамбула-стиль читлистов КР-1 (дискретная математика).
% Подключается из kr1-cheatsheet-1side.tex и kr1-cheatsheet-2side.tex.
\usepackage{fix-cm}
\usepackage[a4paper,margin=6mm]{geometry}
\usepackage{amsmath,amssymb}
\usepackage{fontspec}
\setmainfont{cmunrm.otf}[BoldFont=cmunbx.otf,ItalicFont=cmunti.otf,BoldItalicFont=cmunbi.otf]
\setsansfont{cmunss.otf}[BoldFont=cmunsx.otf,ItalicFont=cmunsi.otf,BoldItalicFont=cmunso.otf]
\usepackage[russian]{babel}
\usepackage{multicol}
\usepackage{enumitem}
\usepackage{xcolor}
\usepackage{array}

\pagestyle{empty}
\setlength{\parindent}{0pt}
\setlength{\parskip}{0pt}
\setlength{\columnsep}{2.4mm}
\setlength{\columnseprule}{0.25pt}
\setlength{\multicolsep}{0pt}
\raggedcolumns
\sloppy
\usepackage{ragged2e}
\AtBeginDocument{\RaggedRight}
\emergencystretch=2em
\setlength{\abovedisplayskip}{1pt}\setlength{\belowdisplayskip}{1pt}
\setlength{\abovedisplayshortskip}{0pt}\setlength{\belowdisplayshortskip}{1pt}
\setlength{\tabcolsep}{2.2pt}
\renewcommand{\arraystretch}{0.95}
\setlength{\arraycolsep}{1.5pt}
\thinmuskip=2mu \medmuskip=2.5mu plus 1mu minus 1mu \thickmuskip=3mu plus 1mu

% основной кегль 8/8.8 pt (не меньше 7 pt — context/prompts/cheatsheet.md §3)
\newcommand{\sheetsize}{\fontsize{8}{8.8}\selectfont}
\DeclareMathSizes{8}{8}{5.8}{4.6}

\setlist{nosep,leftmargin=*,labelsep=2pt,itemsep=0pt,parsep=0pt,topsep=0pt,partopsep=0pt}
\setlist[itemize]{label=\textbullet}
\setlist[enumerate]{label=\arabic*.}

% Заголовок раздела: номер темы КР по билету, светло-серая плашка на всю колонку.
\newcommand{\sect}[1]{\par\addvspace{1.6pt}\noindent
  {\setlength{\fboxsep}{1pt}\colorbox{black!13}{\parbox{\dimexpr\linewidth-2\fboxsep\relax}{\sffamily\bfseries #1}}}%
  \par\nobreak\vspace{0.5pt}\nobreak}
% Подзаголовок внутри раздела.
\newcommand{\sub}[1]{\par\addvspace{0.4pt}\noindent\textbf{#1}\enspace}
% Единый маркер ловушки во всём листе.
\newcommand{\trap}{{\setlength{\fboxsep}{0.4pt}\setlength{\fboxrule}{0.4pt}\fbox{\textbf{\,!\,}}}\ }
% Обозначения курса
\newcommand{\N}{\mathbb{N}}
\newcommand{\Z}{\mathbb{Z}}
\newcommand{\Q}{\mathbb{Q}}
\newcommand{\nil}{[\,]}
\newcommand{\lh}{\mathrm{lh}}
\newcommand{\app}{\mathrm{app}}
\newcommand{\rev}{\mathrm{rev}}
\newcommand{\Prog}{\mathrm{Prog}}
\newcommand{\Ra}{\Rightarrow}
\newcommand{\LRa}{\Leftrightarrow}
\newcommand{\topline}[1]{\noindent{\sffamily\bfseries Дискретная математика (ПИ ФКН ВШЭ) · КР-1 · 24.10.2026 · #1}\par\vspace{1pt}}
